\begin{abstract}
变分自编码器（VAE；\cite{vae}）是近期提出的一类生成模型，它将自顶向下的生成网络与自底向上的识别网络相结合，后者近似执行后验推断。这类模型通常对后验推断作出较强假设，例如后验分布近似按维度因子化，其参数可以通过对观测值的非线性回归来近似。我们的实证结果表明，VAE 目标可能产生过度简化的表示，无法充分利用网络的全部建模能力。本文提出重要性加权自编码器（IWAE）：其架构与 VAE 相同，但使用由重要性加权导出的更紧的对数似然下界。IWAE 的识别网络利用多个样本近似后验，因此能更灵活地表示不满足 VAE 建模假设的复杂后验。实验表明，IWAE 比 VAE 学到更丰富的潜在空间表示，并在密度估计基准上取得更高的测试对数似然。
\end{abstract}
